\subsection{Connected complex fibres in the regular-evaluation topology}
\label{subsec:ftb-connected-complex-fibres}

For a finite-type complex algebra \(A\), give
\(\operatorname{Hom}_{\mathbb C\text{-alg}}(A,\mathbb C)\)
the topology induced by all evaluations \(x\mapsto x(a)\).
On affine scheme charts this is the topology of the actual
complex points used in the uniformization construction.
We record the source's elementary connectedness proof, since it
is used to make the relative exponential surjective.

\begin{srlemma}[Polynomial growth off a hypersurface]
\label{lem:ftb-polynomial-growth-off-hypersurface}
Let \(0\ne q\in\mathbb C[X_1,\ldots,X_d]\).
If \(f\) is holomorphic on \(q\ne0\) and
\[
                 |f(x)|\le C(1+\|x\|)^m \qquad(q(x)\ne0),
                                                               \tag{CC1}
\]
then \(f\) agrees there with a polynomial.
Also \(q\ne0\) is connected and dense in \(\mathbb C^d\).
\end{srlemma}
\begin{proof}
For the last assertion, restrict \(q\) to the complex affine
line joining two points where it is nonzero.
The restriction is a nonzero one-variable polynomial.
The complement of its finite zero set in \(\mathbb C\)
is path connected, so a path from parameter \(0\) to \(1\)
maps to the required path in \(q\ne0\).
For density, join an arbitrary point \(x\) to a point \(b\)
where \(q(b)\ne0\).
The same restriction has only finitely many zeros.
Nonzero parameters tending to \(0\), avoiding those zeros,
give points of \(q\ne0\) tending to \(x\).

For the polynomial assertion in one variable, remove the
roots of \(q\) one at a time.
Bound (CC1) bounds \(f\) near each deleted root, so the
removable-singularity theorem extends \(f\) there.
The polynomial bound persists by continuity.
The extended entire function is a polynomial of degree at
most \(m\), by Cauchy's estimates on large circles.

Proceed by induction on \(d\).
Write the first coordinate as \(t\) and the remaining
coordinates as \(x'\).
There are only finitely many \(t\) for which \(q(t,x')\)
is the zero polynomial in \(x'\): regard \(q\) as a nonzero
polynomial over \(\mathbb C[X']\), and use its finite root set.
Choose \(m+1\) distinct good values \(t_0,\ldots,t_m\).
Induction gives polynomials \(P_j(x')\) agreeing with
\(f(t_j,x')\) wherever \(q(t_j,x')\ne0\).
Form the polynomial
\[
                 P(t,x')=\sum_{j=0}^m
                       \ell_j(t)P_j(x'),                    \tag{CC2}
\]
where \(\ell_j\) are the Lagrange basis polynomials.
For a fixed \(x'\) for which all \(q(t_j,x')\ne0\), the
one-variable assertion applied to \(t\mapsto f(t,x')\)
gives a polynomial of degree at most \(m\).
Its \(m+1\) specified values give (CC2).
The bound needed for both inductions follows from
\[
       1+\|(t,x')\|\le(1+|t|)(1+\|x'\|).
\]
Finally the nonzero polynomial \(\prod_jq(t_j,x')\)
has dense nonvanishing set by the first paragraph.
For any point of \(q\ne0\), perturb \(x'\) toward this set
along a complex line, keeping \(t\) fixed.
The perturbation stays in \(q\ne0\) near that point.
Continuity of \(f\) and \(P\) gives their equality at the
original point as well.
\end{proof}

\begin{srlemma}[The connected unramified root cover]
\label{lem:ftb-connected-unramified-root-cover}
Let \(p\in\mathbb C[X_1,\ldots,X_d][T]\) be monic and
irreducible, and let \(q\ne0\) be a polynomial such that
\(p_x\) has only simple roots whenever \(q(x)\ne0\).
Then
\[
                   Z=\{(x,t):q(x)\ne0,\ p_x(t)=0\}
                                                               \tag{CC3}
\]
is connected.
\end{srlemma}
\begin{proof}
Suppose two nonempty disjoint relatively open sets partition \(Z\).
For each \(x\), form the monic product of \(T-t\) over the
roots lying in the first set, and the analogous complementary
product.
Simple roots have local analytic continuations, by the
one-variable derivative and analytic inverse theorem.
Their membership in either open part is locally constant.
Thus all coefficients of these two partial products are
holomorphic on \(q\ne0\).

The roots of a monic polynomial satisfy
\[
                   |t|\le 1+\sum_{j<\deg p}|p_j(x)|.
\]
Each root therefore has a polynomial bound in \(\|x\|\).
The elementary-symmetric expressions for coefficients of
a product of at most \(\deg p\) such roots have a common
polynomial bound as well.
Lemma~\ref{lem:ftb-polynomial-growth-off-hypersurface}
makes each coefficient a polynomial in \(x\).
Retain coefficients only through degree \(\deg p\); higher
ones are zero in every partial product.
We obtain \(A,B\in\mathbb C[X][T]\) with
\(A_xB_x=p_x\) for \(q(x)\ne0\).
Density and equality of polynomial evaluations give \(AB=p\).
A root in each of the two nonempty parts makes both \(A\)
and \(B\) nonunits: its specialized partial product has that root,
whereas a unit polynomial over \(\mathbb C[X]\) specializes to
a nonzero constant.
This contradicts irreducibility of \(p\).
The set (CC3) is nonempty, since the monic irreducible \(p\)
has positive degree and every complex specialization has a root.
\end{proof}

\begin{srlemma}[Finite integral maps and dense principal opens]
\label{lem:ftb-complex-principal-open-dense}
If \(A\) is a finite-type complex domain and \(0\ne f\in A\),
the points \(x\) with \(x(f)\ne0\) are dense in its complex points.
\end{srlemma}
\begin{proof}
We spell out the finite-map fact used by the source.
If \(R\hookrightarrow A\) is finite, with both rings domains
and \(R\) integrally closed, restriction of complex points is
an open map.
First, restriction for an integral ring map is proper in
the evaluation topology.
For an ultrafilter of points whose restrictions converge,
a monic integral equation bounds every coordinate \(x(a)\)
by \(1+\sum|x(c_j)|\).
The coefficients \(c_j\) on the base are bounded along
the ultrafilter.
Compact complex closed balls therefore supply an
ultrafilter limit for each coordinate.
The coordinate limits preserve addition, multiplication,
constants and \(1\); they give a complex algebra homomorphism
and the required limit point.
This proves properness, and in particular closedness.

To prove openness, let \(x\) be a point in an open set \(U\).
There are only finitely many points in its restriction fibre.
Choose \(a\in A\) with \(x(a)=0\) and \(x'(a)=1\) for
all points \(x'\) of that fibre outside \(U\).
Such an \(a\) is constructed from pairwise separating
functions: if \(e_{x'}(x)=1\), \(e_{x'}(x')=0\), take
\(a=1-\prod e_{x'}\).
The image of
\[
                        U^c\cap\{x':|x'(a)|\le1/2\}
\]
is closed, by properness, and misses the given base point.
Let \(P\in R[T]\) be the monic minimal polynomial of \(a\)
over \(\operatorname{Frac}R\); its coefficients lie in \(R\)
by integral closedness.
At the given base point \(P(0)=0\).
At every sufficiently near base point the absolute value
of its constant coefficient is less than \(2^{-\deg P}\).
The product formula for the roots of a monic polynomial
then gives a root \(t\) with \(|t|<1/2\).
Evaluation \(a\mapsto t\) defines a point of \(R[a]\):
the kernel of \(R[T]\to R[a]\) is \((P)\), by the
minimal-polynomial divisibility theorem over an integrally
closed domain.
It extends to a point of \(A\) by integral lying over and
algebraic closedness of \(\mathbb C\).
A base point also outside the preceding closed image thus
has a lift in \(U\).
These two neighborhood conditions prove openness.

Now choose Noether normalization
\(R=\mathbb C[X_1,\ldots,X_d]\hookrightarrow A\), finite and
injective.
Integrality of \(f\ne0\) gives \(0\ne b\in R\cap(f)\):
use the nonzero constant coefficient of a least-degree
monic equation, or the corresponding algebraic divisibility
statement.
The image of the zero set of \(f\) lies in \(b=0\).
This hypersurface has real Lebesgue measure zero, by induction
on \(d\) and Fubini (a nonzero one-variable polynomial has
finitely many roots).
Its complement is dense.
For any nonempty open set of points of \(A\), openness of
restriction gives a nonempty open set in \(\mathbb C^d\).
Choose its base point off \(b=0\); any lift in the given
open set has \(f\ne0\).
This proves density.
\end{proof}

\begin{srproposition}[Connectedness of an integral complex fibre]
\label{prop:ftb-integral-complex-points-connected}
The complex points of an integral locally finite-type
complex scheme are connected in the regular-evaluation
topology.
Consequently, in the complex points of a smooth
geometrically connected algebraic group, a subgroup
containing a neighborhood of the identity described by
finitely many regular-function tests is the whole group.
\end{srproposition}
\begin{proof}
First let the scheme be \(\operatorname{Spec}A\), with
\(A\) a finite-type complex domain.
The source refines Noether normalization to the following
explicit primitive-element data:
\[
 R=\mathbb C[X_1,\ldots,X_d]\hookrightarrow A,\quad a\in A,
 \quad p=\min_R(a),\quad0\ne q\in R,                       \tag{CC4}
\]
where \(A\) is finite over \(R\), \(p\) is monic irreducible,
\[
 q\in(p,p')\subset R[T],\qquad
             qA\subset R[a].                            \tag{CC5}
\]
For completeness, the finite fraction-field extension is
separable in characteristic zero.
A primitive element can be multiplied by a base
denominator to lie in \(A\).
Write finitely many \(R\)-module generators of \(A\) as
rational polynomials in that element; their common
denominator gives the second condition of (CC5).
A Bezout relation between the fraction-field polynomials
\(p,p'\), with denominators cleared, gives the first.
Multiplying those two nonzero denominators supplies one \(q\).
Integral closedness of \(R\) identifies the kernel of
\(R[T]\to A\) with \((p)\).

The conditions (CC5) identify the points on \(q\ne0\)
with the root cover \(q(x)\ne0,\ p_x(t)=0\).
Indeed choose \(r_b(T)\in R[T]\) with \(qb=r_b(a)\);
the inverse point is
\[
                            b\longmapsto r_b(x,t)/q(x).
                                                               \tag{CC6}
\]
Minimal-polynomial divisibility shows that (CC6)
preserves sums and products and is independent of the
chosen \(r_b\).
It is continuous in every evaluation coordinate.
The Bezout relation in (CC5) makes these roots simple.
Lemma~\ref{lem:ftb-connected-unramified-root-cover}
therefore gives a connected principal open.

More generally, for any nonzero \(f\in A\), choose
\(0\ne q'\in R\) divisible by \(f\) in \(A\), say \(q'=ff'\).
The same construction with base nonvanishing \(qq'\)
gives a connected principal open \(D(qf'f)\) inside \(D(f)\).
It is dense in all points, by
Lemma~\ref{lem:ftb-complex-principal-open-dense}.
Hence \(D(f)\), lying between that connected set and its
closure, is connected.
Taking \(f=1\) proves the affine assertion.

For an integral locally finite-type scheme, two complex
points lie in affine opens \(U,V\).
These opens have nonempty intersection by irreducibility,
and that nonempty finite-type open contains a complex
point by the weak Nullstellensatz.
The connected sets \(U(\mathbb C),V(\mathbb C)\) therefore
meet and their union connects the two points.
This proves connectedness in general.

For the last assertion take the actual complex fibre of
the algebraic group.
Smoothness and geometric connectedness over the
algebraically closed field make that fibre integral.
For a subgroup \(H\) containing the given regular-test
neighborhood of the identity, membership in \(H\) is
locally constant: translate that neighborhood by any
fixed point; points in the resulting neighborhood have
the same membership as the fixed point.
Translations are continuous because they are algebraic
and evaluation is compatible with pullback.
Connectedness makes this membership constant.
It holds at the identity, so \(H\) is the whole group.
\end{proof}

The same affine argument gives a useful finite-type variant:
any locally constant function on the complex points of
a finite-type complex scheme has finite image.
There are finitely many minimal-prime components on an
affine chart; the function is constant on each corresponding
integral complex point space.
A finite affine cover gives finite image globally.
Nilpotents do not affect complex points.

\paragraph{Frozen source concordance.}
All paths in the first list have prefix
\srcid{src/Development/AbelianVarietyTheory/Proofs/ToricGeometry/}.
The hypersurface argument is
\srcid{IsConnectedAndDenseComplHypersurface.lean}
(\srcid{c7a30e8c1c08de636533db8f1d17463a42588382}).
The polynomial extension, using removable singularities
and the finite interpolation (CC2), is
\srcid{ExistsMvPolynomialOfPolyBoundOffHypersurface.lean}
(\srcid{b5df170413a5d462c24189a4132bc4738bb9e632}).
The partial-root products are
\srcid{IsConnectedUnramifiedCoverOfIrreducible.lean}
(\srcid{19755400e2e92e979664c3b03ed114a0942d13ae}).
The primitive-element and common-denominator construction is
\srcid{ExistsPrimitiveElementDiscriminantMultiple.lean}
(\srcid{bf8c8d866c3dfba43b74d7168b6a47d61fbb8212}).
The actual point formula (CC6) and principal-open reduction
are
\srcid{ExistsIsConnectedPrincipalOpen.lean}
(\srcid{f607bc0b550282ed92207e0dc3e22e3af3f6c25b}),
\srcid{PrincipalOpenDenseComplexPoints.lean}
(\srcid{b86e5fc2d14e1b6c6a7fe6b68c304462a9399278}),
and
\srcid{IsConnectedBasicOpenComplexPointsOfDense.lean}
(\srcid{3e09e88d4274d9fd562649d86140f267a2d55ed1}).
The proper/open integral-map argument is
\srcid{src/Development/AbelianVarietyTheory/Proofs/HomogeneousIdeals/ExistsApplyEqZeroAndEqOne.lean}
(\srcid{ac0ae26397e0083ee3c61fa6d1bd0cbac52acb56});
its density application is
\srcid{ExistsNeZeroIsStandardSmoothOfRelativeDimension.lean}
(\srcid{8272e270db53df9168851be7f7b8da698651807d})
in that same directory.
The passage to actual affine scheme points, minimal-prime
components and finite image is
\srcid{src/Development/AbelianVarietyTheory/Proofs/ModuliProofs/FiniteImageOfLocallyConstantOnHolomorphicFamilie.lean}
(\srcid{6cd65b6f6d687b4dde3adf5384e111ee945bd6bf}).
The general integral-fibre and subgroup argument is
\srcid{PreconnectedSpacePtsOfIsIntegral.lean}
(\srcid{1189a01139efe76d2e9228654d277183847db379})
in that directory.
The argument proves the connectedness input of the
relative exponential; it assumes no PEL uniformization,
no family of theta frames, and no algebraization of an
arbitrary complex torus.
